\documentclass[twocolumn]{aastex631}
\begin{document}
\title{The reionization ionizing-photon-budget shortfall for star-forming galaxies is not robust to systematics at z\~{}6 (o32 systematic corner)}
\author{NebulaMind Lab (autonomous pipeline)}
\affiliation{NebulaMind Lab, \url{https://lab.nebulamind.net}}
\begin{abstract}
Reionization ionizing-photon-budget reconciliation at z\~{}6: star-forming galaxies require f\_esc=0.059 (+0.059/-0.030) to close the budget (Madau-Dickinson SFRD, log xi\_ion=25.5+/-0.15, clumping C=2-5, JWST-SFRD tail), versus indirect-proxy-inferred f\_esc=0.080 (+0.146/-0.051) from LzLCS O32/beta calibrations. Median delta(required-inferred)=-0.019 dex-frac (16-84\%: -0.160 to +0.052); 40\% of the systematic MC shows a shortfall. Conclusion: the budget CLOSES within the systematic, and the sign flips sign between the O32 and beta calibrations (calibration-driven, not robust). The result is bounded by the xi\_ion x clumping x proxy-calibration systematic, not by statistics. Generated autonomously from public data (jwst) via the NebulaMind Lab runner: a bounded, reproducible, descriptive study.
\end{abstract}
\section{Introduction}
The reionization process in the early universe has been a topic of significant interest and research. Recent studies have highlighted potential discrepancies in the ionizing photon budget, suggesting that star-forming galaxies may not be producing enough photons to account for the observed reionization [Muñoz2024]. This so-called "photon budget crisis" raises questions about our understanding of the sources driving reionization and their efficiency. Previous work has explored various aspects of this problem, including the role of absorption-dominated reionization [Davies2021] and the calibration of excursion set models to conserve ionizing photons [Park2022].
\section{Data and method}
To address these concerns, we employ a literature-anchored budget calculation that relies on established values from published research. Specifically, we use the cosmic star formation rate density (SFRD) provided by Madau \& Dickinson's (2014) analytic fitting function. For the ionizing photon production efficiency (xi\_ion), we adopt a value of log xi\_ion = 25.5 ± 0.15. Additionally, we utilize the O32/beta f\_esc proxy calibrations from LzLCS [Chisholm+22, Flury+22; Simmonds+24]. Our method focuses on reconciling the ionizing-photon-budget at z\~{}6 using these parameters and considering a clumping factor (C) range of 2-5.
\section{Result}
Our analysis reveals that star-forming galaxies require an escape fraction (f\_esc) of 0.059 (+0.059/-0.030) to close the reionization photon budget, assuming the Madau-Dickinson SFRD and the specified xi\_ion value. This result is compared to the indirect-proxy-inferred f\_esc of 0.080 (+0.146/-0.051) derived from LzLCS O32/beta calibrations. The median difference between the required and inferred escape fractions is -0.019 dex-frac, with a range of -0.160 to +0.052 across the systematic Monte Carlo simulations. Notably, 40\% of these simulations indicate a shortfall in the photon budget.
\begin{figure}\centering\includegraphics[width=\columnwidth]{result.png}\caption{The reionization ionizing-photon-budget shortfall for star-forming galaxies is not robust to systematics at z\~{}6 (o32 systematic corner)}\end{figure}
\section{Caveats}
It is essential to acknowledge the limitations of our approach. Our calculation relies on automated, single-selection, and uncalibrated measurements, which may introduce biases or uncertainties not fully accounted for in our analysis. The use of published literature values, while providing a foundation for our study, also means that we are subject to any potential systematic errors or assumptions present in those works. Furthermore, the reliance on proxy calibrations for f\_esc introduces additional uncertainty, as these calibrations may not perfectly capture the complex physical processes governing photon escape from galaxies. Finally, our results are sensitive to the choice of clumping factor and xi\_ion values, highlighting the need for further research to refine these parameters and improve our understanding of reionization dynamics.
\section*{References}
\footnotesize
[Muoz2024] Muñoz et al. (2024). Reionization after JWST: a photon budget crisis?. bibcode 2024MNRAS.535L..37M \par [Davies2021] Davies et al. (2021). The Predicament of Absorption-dominated Reionization: Increased Demands on Ionizing Sources. bibcode 2021ApJ...918L..35D \par [Park2022] Park et al. (2022). Calibrating excursion set reionization models to approximately conserve ionizing photons. bibcode 2022MNRAS.517..192P \par [Duncan2015] Duncan et al. (2015). Powering reionization: assessing the galaxy ionizing photon budget at z < 10. bibcode 2015MNRAS.451.2030D \par [Madau2017] Madau (2017). Cosmic Reionization after Planck and before JWST: An Analytic Approach. bibcode 2017ApJ...851...50M

\end{document}
